\chapter{API Reference}

La documentazione interattiva è disponibile automaticamente all'avvio del server
su \url{http://localhost:8000/docs} (Swagger UI) e \url{http://localhost:8000/redoc}.

\section{Health Check}

\begin{table}[H]
  \centering
  \begin{tabularx}{\textwidth}{lX}
    \toprule
    \textbf{Campo} & \textbf{Valore} \\
    \midrule
    Metodo & \texttt{GET} \\
    Path   & \texttt{/health} \\
    Descrizione & Verifica che il server sia attivo \\
    \bottomrule
  \end{tabularx}
\end{table}

\textbf{Risposta 200:}
\begin{lstlisting}[language=Python]
{"status": "ok", "version": "1.0.0"}
\end{lstlisting}

\section{Endpoint Documenti}

\subsection{Carica un documento}

\begin{table}[H]
  \centering
  \begin{tabularx}{\textwidth}{lX}
    \toprule
    \textbf{Campo} & \textbf{Valore} \\
    \midrule
    Metodo      & \texttt{POST} \\
    Path        & \texttt{/documents/upload} \\
    Content-Type & \texttt{multipart/form-data} \\
    Parametro   & \texttt{file} (UploadFile, required) \\
    \bottomrule
  \end{tabularx}
\end{table}

\textbf{Esempio richiesta:}
\begin{lstlisting}[language=bash]
curl -X POST http://localhost:8000/documents/upload \
  -F "file=@/path/to/document.pdf"
\end{lstlisting}

\textbf{Risposta 201:}
\begin{lstlisting}[language=Python]
{
  "id": "3fa85f64-5717-4562-b3fc-2c963f66afa6",
  "filename": "documento_abc123.pdf",
  "upload_date": "2026-06-11T10:30:00+00:00",
  "total_chunks": 47,
  "total_pages": 8,
  "file_size_bytes": 524288,
  "status": "ready"
}
\end{lstlisting}

\textbf{Errori:}
\begin{itemize}
  \item \texttt{400} — File non PDF, file vuoto, o tipo MIME non valido
  \item \texttt{413} — Dimensione superiore al limite configurato (default 50 MB)
  \item \texttt{500} — Errore durante l'indicizzazione
\end{itemize}

\subsection{Lista documenti}

\begin{lstlisting}[language=bash]
curl http://localhost:8000/documents/
\end{lstlisting}

\textbf{Risposta 200:} Array di oggetti documento, ordinati per data di upload decrescente.

\subsection{Dettaglio documento}

\begin{lstlisting}[language=bash]
curl http://localhost:8000/documents/{doc_id}
\end{lstlisting}

\subsection{Elimina documento}

\begin{lstlisting}[language=bash]
curl -X DELETE http://localhost:8000/documents/{doc_id}
\end{lstlisting}

Elimina il documento da PostgreSQL, i vettori da Qdrant e il file dal disco.
Risposta: \texttt{204 No Content}.

\section{Endpoint Chat (RAG)}

\subsection{Fai una domanda}

\begin{table}[H]
  \centering
  \begin{tabularx}{\textwidth}{lX}
    \toprule
    \textbf{Campo} & \textbf{Valore} \\
    \midrule
    Metodo      & \texttt{POST} \\
    Path        & \texttt{/chat/} \\
    Content-Type & \texttt{application/json} \\
    \bottomrule
  \end{tabularx}
\end{table}

\textbf{Body JSON:}
\begin{lstlisting}[language=Python]
{
  "question": "Qual e la metodologia descritta nel capitolo 3?",
  "doc_id": "3fa85f64-...",   // opzionale: filtra su un documento
  "session_id": "abc-123",    // opzionale: per mantenere lo storico
  "top_k": 4,                 // default: 4, range: 1-20
  "stream": false             // true per SSE streaming
}
\end{lstlisting}

\textbf{Risposta 200 (non-streaming):}
\begin{lstlisting}[language=Python]
{
  "session_id": "abc-123",
  "answer": "Il capitolo 3 descrive una metodologia basata su...",
  "sources": [
    {
      "text": "...testo del chunk rilevante...",
      "page_num": 12,
      "doc_id": "3fa85f64-...",
      "chunk_index": 23,
      "score": 0.8932
    }
  ]
}
\end{lstlisting}

\textbf{Risposta streaming (\texttt{"stream": true}):}\\
Server-Sent Events (SSE). Ogni evento contiene un token della risposta.
L'ultimo evento è \texttt{data: [DONE]}.

\begin{lstlisting}[language=bash]
curl -X POST http://localhost:8000/chat/ \
  -H "Content-Type: application/json" \
  -d '{"question": "Di cosa tratta?", "stream": true}'
# Output:
# data: Il
# data:  documento
# data:  tratta
# data:  di...
# data: [DONE]
\end{lstlisting}

\subsection{Storico conversazione}

\begin{lstlisting}[language=bash]
curl http://localhost:8000/chat/history/{session_id}
\end{lstlisting}

\textbf{Risposta 200:} Array di messaggi con \texttt{role} (\texttt{user}/\texttt{assistant})
e \texttt{content}.

\section{Endpoint Ricerca Semantica}

\begin{table}[H]
  \centering
  \begin{tabularx}{\textwidth}{llX}
    \toprule
    \textbf{Parametro} & \textbf{Tipo} & \textbf{Descrizione} \\
    \midrule
    \texttt{q}       & string (required) & Testo da cercare semanticamente \\
    \texttt{top\_k}  & integer (1-20)    & Numero di risultati, default 5 \\
    \texttt{doc\_id} & string (optional) & Filtra su un singolo documento \\
    \bottomrule
  \end{tabularx}
\end{table}

\textbf{Esempio:}
\begin{lstlisting}[language=bash]
curl "http://localhost:8000/search/?q=reti+neurali&top_k=3&doc_id=3fa85f64-..."
\end{lstlisting}

\textbf{Risposta 200:}
\begin{lstlisting}[language=Python]
{
  "query": "reti neurali",
  "doc_id": "3fa85f64-...",
  "total": 3,
  "results": [
    {
      "text": "Le reti neurali profonde sono composte da...",
      "page_num": 7,
      "doc_id": "3fa85f64-...",
      "chunk_index": 15,
      "score": 0.9214
    }
  ]
}
\end{lstlisting}

\section{Risultati Sperimentali}

Di seguito tre domande di test eseguite sul documento \texttt{sample.pdf} incluso
nella cartella \texttt{tests/fixtures/}. Le risposte dimostrano il corretto funzionamento
della pipeline RAG.

\subsection{Test 1: Domanda Generale}

\textbf{Domanda:} \emph{"Di cosa tratta questo documento?"}

\textbf{Risposta attesa:}
Il documento è un documento di esempio per testare il sistema RAG. Descrive un
AI Document Assistant capace di estrarre testo da PDF, suddividerlo in chunk e
indicizzare gli embedding in Qdrant.

\subsection{Test 2: Architettura}

\textbf{Domanda:} \emph{"Quali sono i componenti principali dell'architettura?"}

\textbf{Risposta attesa:}
L'architettura include FastAPI come framework web, LangChain e LangGraph per
l'orchestrazione RAG, Qdrant come database vettoriale e PostgreSQL per i metadati.

\subsection{Test 3: Deployment}

\textbf{Domanda:} \emph{"Come si avvia il progetto?"}

\textbf{Risposta attesa:}
Il progetto si avvia con il comando \texttt{docker compose up}. La Swagger UI
è disponibile su \texttt{http://localhost:8000/docs}.
